%% 第三章 图增强检索与结构化问答

\chapter{图增强检索与结构化问答}
\chaptermark{图增强检索与结构化问答}

\section{引言}

如第二章所述，层级情景记忆树通过将原始多模态经历组织为 L0–L4+ 的多层纵向结构，配合大语言模型 Agent 的交互式树导航，以远低于"将完整历史直接拼接输入模型"的上下文成本实现了长历史问答。然而，树的检索范式是纯纵向的，Agent 沿父子关系展开节点，证据召回的上限由树的层级结构和文本语义相似度决定。许多回溯性问题天然依赖跨任务、跨时间的横向联系：例如"这个杯子后来出现在哪些任务中""那次失败之前发生了什么""同一位置还有哪些事件"，这些问题的答案分散在树的不同分支和层级中，仅靠纵向展开和文本搜索极易遗漏关键证据。

本章解决第一章所提出的第一个研究问题（RQ1）：如何在长历史条件下，弥补纯层级结构和文本相似度检索在横向语义关联上的不足。核心方法是在不修改树存储结构的前提下，在树结构之上叠加一层\textbf{事件关联图}索引：从 L2 事件节点自动派生图节点，构建时间、物体、位置、动作和因果五类语义边，将横向联想能力嵌入检索链路。同时，将记忆树中的任务、物体、日期等结构化信息显式化为可调用工具，并配合预处理缓存和执行护栏，使 Agent 在长历史上的搜索行为从"无结构的自由展开"转变为"沿结构化路径精确命中"。

本章其余部分安排如下。3.2 节介绍事件关联图的构建方法，包括节点派生、五类边的定义与构建算法、密度控制与缓存策略。3.3 节介绍图增强检索的三阶段流水线，包括语义种子召回、图邻居扩展与联合重排。3.4 节介绍结构化检索工具与执行护栏的设计。3.5 节在 TEACh 数据集上报告两种长历史设置下的问答实验结果，并对主要机制进行归因分析。

\begin{figure}[htbp]
    \centering
    \includegraphics[width=\textwidth]{图3-1 图增强检索系统整体架构图.png}
    \caption{图增强检索系统整体架构图}
    \label{fig:chapter3-architecture}
\end{figure}

\section{事件关联图构建}

事件关联图是 GRAF-Mem 的核心数据结构。它从层级记忆树自动派生，作为树之上的横向索引层，使检索同时具备"沿树上下搜索"和"沿图横向扩展"两条路径。更重要的是，图结构在后续章节中还被遗忘模块用于中心性评估（第四章）和修正模块用于传播检测（第五章）。这种"一图多用"的设计是 GRAF-Mem 的核心创新之一。本节详细阐述图节点与边的定义和构建算法。

\subsection{图节点的定义与派生}

事件关联图的节点直接对应层级记忆树中的 L2 层节点（L2事件摘要节点）。选择 L2 层作为图节点来源的原因有三：（1）L2 节点是记忆树中最细粒度的自然语言层，每个 L2 节点对应一个具体的操作动作（如 "Pickup(Mug\_0)"）或一条用户语音转写，颗粒度适合作为检索的基本单元；（2）L2 节点携带的信息维度最为丰富，包含物体集合（从 L1 场景图汇总而来）、动作标签、位置签名和时间范围，可以支撑多种不同类型的边构建；（3）L2 节点的数量适中，在由 50 个基础情景拼接得到的长历史设置下约有 5000–6000 个节点，构建和缓存图的成本可控。

具体派生流程如下。设记忆树共有 $N$ 个 L2 节点 $\{E_1, E_2, \ldots, E_N\}$，对每个节点 $E_i$ 提取以下属性：

\begin{itemize}
\item \textbf{物体实例 ID 集合} $O(E_i)$：遍历 $E_i$ 关联的所有 L1 场景图，汇总各帧中出现的物体实例 ID（如 "Mug\_0"、"CounterTop\_1"）
\item \textbf{主要动作} $A(E_i)$：从 $E_i.\mathrm{latest\_raw}.\mathrm{current\_action}$ 提取，如 "Pickup(Mug\_0)"、"Place(CounterTop\_0)"、"Noop"
\item \textbf{位置签名集合} $P(E_i)$：由于 TEACh 数据没有显式的房间标签，位置签名通过各帧场景图中的大型固定物体（家具、电器）的类别集合来近似。具体地，取该事件最后一个场景图中属于位置指示物体集合（如 Sink、Fridge、CounterTop、Cabinet、StoveBurner、Toaster、Sofa、Bed 等 20 余类家具/电器）的物体类别集合，排序后作为位置签名集合
\item \textbf{时间范围} $T(E_i) = [t_{\mathrm{start}}, t_{\mathrm{end}}]$：从该事件对应的时间范围提取
\end{itemize}

每个图节点 $G_i$ 继承其源 L2 事件 $E_i$ 的上述属性，且保留对原始树节点 $E_i$ 的引用（$G_i.\mathrm{tree\_node} = E_i$），以支持从图节点回树节点进行细节展开。

\subsection{五种边类型的定义与构建算法}

事件关联图支持五种边类型，分别编码不同类型的事件间横向语义关联。每种边类型的构建算法独立运行，最终的图为它们的并集。以下按构建成本从低到高逐一阐述。

\textbf{（1）时间相邻边}

时间相邻边连接同一 L3 目标摘要节点（Goal）下前后相邻的两个 L2 事件。其构建为纯规则驱动，零嵌入计算成本。构建算法为：对每个 L3 目标摘要节点 $g$，取其所包含的 L2 事件序列 $[E_1, E_2, \ldots, E_k]$，对序列中每对相邻事件 $(E_i, E_{i+1})$ 添加一条时间相邻边，并记录前后方向标签；检索时可根据问题中的方向词选择前向或后向展开。边的权重为固定的 1.0（因为时间相邻是确定性的结构关系，不存在"强度"变化）。

该边类型支撑"之前/之后发生了什么"类时序推理问题。例如，用户问 "what did you do just after making coffee"，Agent 可以从"making coffee"对应的图节点出发，沿时间相邻边定位到紧后的事件并展开其摘要。

\textbf{（2）共享物体边}

共享物体边连接涉及相同物体实例 ID 的两个 L2 事件。构建算法为：对每一对共享至少一个物体实例 ID 的事件 $(E_i, E_j)$，添加一条无向边。边的权重为两个事件物体集合的 Jaccard 系数：

\begin{equation}
w_{\mathrm{co\_object}}(E_i, E_j) = \frac{|O(E_i) \cap O(E_j)|}{|O(E_i) \cup O(E_j)|}
\label{eq:3-1}
\end{equation}

为避免过于常见的物体（如 "agent hand"、"Floor" 等几乎出现在所有事件中的物体）导致图过密，对纳入 Jaccard 计算的物体集合进行过滤：排除全局频率超过阈值（默认为所在 Goal 的事件数的 80\%）的物体 ID。

该边类型支撑"与某个物体相关的所有事件"类问题。例如，用户问 "where else did the blue mug appear"，Agent 可以从当前已知包含蓝色杯子的图节点出发，沿共享物体边找到该杯子出现的所有其他事件位置。

\textbf{（3）共位置边}

共位置边连接共享相同位置签名（即场景中出现的家具/电器集合高度重叠）的两个 L2 事件。构建算法为：对每对事件计算其位置签名集合的 Jaccard 系数（基于位置物体类别集合而非实例 ID），若 $\ge$ 阈值（默认 0.5），则添加一条无向边。边的权重即为位置签名的 Jaccard 系数。

该边类型支撑"同一位置还发生了哪些事"类空间推理问题。例如，用户问 "what else happened in the kitchen"，Agent 可以从 kitchen 相关事件出发，沿共位置边找到同一空间内的其他事件。

\textbf{（4）相似动作边}

相似动作边连接动作语义相近的 L2 事件。动作相似度通过 Sentence-BERT 嵌入模型（all-mpnet-base-v2）计算。直接比较所有事件对的动作文本会带来很高的计算开销。例如，当图中有约 6000 个事件时，逐一比较任意两个事件需要进行数千万次相似度计算。为降低开销，本文先对动作文本进行模板化归一，将包含具体对象编号的动作转化为更通用的动作模板。例如，"Pickup(Mug\_0)" 和 "Pickup(Apple\_1)" 都可归一为 "Pickup(object)"，"Place(CounterTop\_0)" 可归一为 "Place(surface)"。这样，大量事件会共享同一个动作模板。

在得到动作模板后，系统只需要比较不同动作模板之间的语义相似度，而不必直接比较所有事件对。设图中共有 $N$ 个事件，归一后得到 $M$ 个不同的动作模板，通常 $M \ll N$，因此计算规模由"事件两两比较"降低为"动作模板两两比较"。最后，将模板之间的相似关系映射回具体事件：若两个动作模板的余弦相似度超过阈值（默认 0.85），则涉及这两个模板的事件之间添加相似动作边，边权重等于动作模板之间的余弦相似度。

该边类型支撑"类似的操作还执行过哪些"类泛化检索。例如，用户问 "when else did you clean something"，Agent 可以找到所有包含 "clean" 相关动作（如 "rinse"、"wipe"、"wash"）的事件。

\textbf{（5）因果边}

因果边尝试连接存在因果关联或失败传播关系的两个事件。其构建依赖大语言模型推断，成本较高，因此默认为可选启用，并不纳入主实验的默认图配置。构建算法为：对每个 Goal 内的事件序列 $[E_1, \ldots, E_k]$，识别其中的异常事件（动作状态为 failure 或目标状态发生突变），将异常事件的上下文（前后各 2 个事件的摘要）发送给大语言模型，由其判断任意两个事件间是否存在因果关联（如"因为上一步把杯子掉到地上打碎了，所以下一步清扫地面"）。若判断存在因果关系，则添加一条有向因果边，权重由模型给出（通常为 0.7–0.9）。

\subsection{密度控制、内容签名与图缓存}

五种边类型中，除时间相邻边的数量随事件数线性增长外，其余四种如果不加控制，都可能产生接近"任意两个事件相连"的稠密图结构。实际构建中，通过三重密度控制确保图的规模可控：（1）\textbf{全局频率过滤}：对共享物体边，过滤掉在超过 80\% 事件中都出现的物体 ID（如 "Floor"、"agent hand"）；（2）\textbf{节点度上限}：对每个节点，每种边类型最多保留 top-k（默认 20）条权重最高的边；（3）\textbf{按边类型权重排序}：在有节点度上限约束的情况下，对候选边按权重降序选择，优先保留关联强度最高的边。

评测管线中，一段待评测的长历史在程序中表示为 \texttt{history} 对象。这里的 \texttt{history} 由一组基础情景（episode）按时间顺序拼接得到，因此 episode 列表可以唯一确定该长历史的内容。对于同一段长历史的多次评测调用，系统共享同一图结构。图的缓存通过历史结构签名实现，具体做法是将节点类型、节点数量、时间范围以及摘要文本内容共同纳入确定性哈希，作为缓存键。若缓存命中，则直接加载已构建的图对象，避免重复的嵌入计算和边构建。在由 50 个基础情景拼接得到的长历史设置下，图通常包含约 5000–6000 个节点和 30000–50000 条边，构建时间约 10–20 秒（包括动作嵌入计算），缓存命中后加载仅需约 1 秒。

图构建在记忆巩固（遗忘）之后执行。遗忘模块先压缩低效用节点，随后在保留节点上构建图；节点数本身并不减少，但部分节点的物体、关系和辅助字段会被清空，因而可用于连边的有效信息减少，最终表现为边更稀疏、检索更聚焦。这种"先筛后连"的管线设计是第三章和第四章模块之间最重要的协作接口。

\begin{figure}[htbp]
    \centering
    \includegraphics[width=\textwidth]{图3-2 事件关联图构建流程.png}
    \caption{事件关联图构建流程}
    \label{fig:chapter3-graph-construction}
\end{figure}

\section{图增强检索流水线}

图增强检索在 Agent 执行 \texttt{history.search(query)} 时被触发，将原本的纯文本相似度 top-k 返回替换为经过图扩展和联合重排后的结果列表。整个流水线包含三个阶段。

\begin{figure}[htbp]
    \centering
    \includegraphics[width=\textwidth]{图3-3 检索流水线图.png}
    \caption{图增强检索三阶段流水线}
    \label{fig:chapter3-pipeline}
\end{figure}

\subsection{阶段一：语义种子召回}

语义种子召回沿用了层级记忆树原有的语义搜索机制。使用 Sentence-BERT 句向量模型 all-mpnet-base-v2（768 维嵌入）计算查询向量 $q$ 和每个节点递归索引内容 $\text{索引内容}(n)$（即节点 $n$ 的全部可检索文本片段）之间的余弦相似度。采用最大值池化策略：$\mathrm{sim\_vec}(q, n) = \max(\mathrm{cos\_sim}(\mathrm{emb}(q), \mathrm{emb}(t)) \text{ for } t \text{ in } \text{索引内容}(n))$，确保只要节点中有一小段文本与查询高度相关即可进入候选集。FAISS flat 索引用于加速大规模检索。

语义种子召回的产出是候选集 $S_0$，包含 top-$k_{\mathrm{sem}}$（默认 $k_{\mathrm{sem}} = 10$）个语义相似度最高的节点。$S_0$ 作为后续图扩展的起点。

\subsection{阶段二：图邻居扩展}

对 $S_0$ 中的每个种子节点 $s$，遍历其通过激活边类型集合 $E_{\mathrm{active}}$ 连接的邻居。$E_{\mathrm{active}}$ 不是固定的，而是由\textbf{查询意图路由器}根据用户问题中的关键词与实体线索动态决定。具体规则可概括为：时间线索词（如 "right after"、"right before"、"before"、"after"、"next"、"previous" 以及中文的"之前/之后"）激活时间相邻边；明确物体名激活共享物体边；空间位置词（如 "where"、"in the kitchen"、"at the sink"）激活共位置边；"why"、"cause"、"because"、"failure" 等因果词激活因果边。若多个条件同时命中，则对应边类型并行激活；若未命中任何规则，则默认激活时间相邻、共享物体、共位置和相似动作边，不启用因果边。

对于每种激活边类型 $e$，从 $s$ 出发最多取 $\mathrm{max\_neighbors}$（默认 10）个权重最高的邻居，所有邻居的集合为扩展集 $S_1$。图邻居扩展的直观意义是：如果节点 $s$ 与用户问题高度语义相关，那么与 $s$ 具有时间/物体/位置/动作关联的邻居节点也可能包含对回答有用的信息；即使这些邻居的文本与问题在措辞层面并不相似，它们仍可能提供关键证据。

\subsection{阶段三：联合重排}

将 $S_0$ 和 $S_1$ 合并为候选池 $S = S_0 \cup S_1$，对 $S$ 中的每个节点 $n$ 计算联合重排分数。记 $\mathrm{neighbor\_sim}(q, n, e)$ 为节点 $n$ 在边类型 $e$ 下所有邻居与查询 $q$ 的语义相似度均值，即 $\mathrm{neighbor\_sim}(q, n, e) = \text{mean}\{ \mathrm{sim\_vec}(q, n') \mid n' \in \mathrm{neighbors}(n, e) \}$。若某一边类型下没有邻居，则该项记为 0。则联合重排分数为：

\begin{equation}
\begin{aligned}
\mathrm{score}(q, n) &= \alpha \cdot \mathrm{sim\_vec}(q, n) \\
  &\quad + \beta \sum_{e \in E(n) \cap E_{\mathrm{active}}} w(e) \cdot \mathrm{neighbor\_sim}(q, n, e)
   + \gamma \cdot \mathrm{depth\_bonus}(n)
\end{aligned}
\label{eq:3-2}
\end{equation}

其中三项含义分别为：

\begin{itemize}
\item \textbf{第一项（$\alpha$ 项）}：节点自身与查询的语义相似度。这是原有纯语义搜索的能力，确保检索结果始终以语义相关性为基底。
\item \textbf{第二项（$\beta$ 项）}：图邻居的加权贡献。对于节点 $n$ 通过激活边类型连接的所有邻居 $n'$，计算邻居与查询的语义相似度，按边权重 $w(e)$ 加权求和后取平均。这一项使那些自身语义相似度一般、但有多个高相关邻居的节点获得提权，因为这些节点可能是连接多个相关事件的结构枢纽。
\item \textbf{第三项（$\gamma$ 项）}：深度奖励。$\mathrm{depth\_bonus}(n) = (\mathrm{max\_depth} - \mathrm{depth}(n)) / \mathrm{max\_depth}$，其中 $\mathrm{max\_depth}$ 为当前树的深度（默认 6），$\mathrm{depth}(n)$ 为节点所在层级（根为 0）。层级越浅，获得越高的深度奖励，因为在记忆树中，浅层节点覆盖的时间范围和语义内容更广，对概览类问题更有价值。
\end{itemize}

其中 $E(n)$ 表示节点 $n$ 实际连接到的边类型集合，$\mathrm{neighbors}(n, e)$ 表示节点 $n$ 在边类型 $e$ 下的邻居集合，$w(e)$ 为该边类型的归一化权重（规则边默认取 1.0，大语言模型（LLM）因果边按推断置信度取值），$\mathrm{neighbor\_sim}(q, n, e)$ 为邻居语义相似度均值。权重默认设置为 $\alpha=0.70$, $\beta=0.25$, $\gamma=0.05$。$\alpha$ 的主导地位（0.70）确保语义相关性始终是排序的最主要依据；$\beta=0.25$ 为图结构信号提供了足够的表达空间，但不会使其压倒语义相关性；$\gamma=0.05$ 给予浅层节点微弱的偏好，在不改变排序主导逻辑的前提下提供结构先验。

最终，按联合重排分数降序返回 top-$k_{\mathrm{rerank}}$ 个节点（与 $S_0$ 的 $k_{\mathrm{sem}}$ 相同，默认 10）给 Agent，Agent 随后可以在这些节点上执行 expand/search/collapse 等进一步操作。

\section{结构化检索工具与执行护栏}

\subsection{工具设计：将树结构暴露为可调用能力}

除了图增强的语义搜索，GRAF-Mem 还将层级记忆树中的结构性信息显式化为七种可调用工具。这些工具直接利用树的时间、任务和物体结构进行确定性或半确定性查找，避免 Agent 在长历史文本中进行无约束搜索。每个工具的内部实现均不涉及额外的大语言模型（LLM）调用；它们直接读取树节点中的结构化字段并格式化返回。

\textbf{（1）task\_list(max\_tasks=30)：}按时间顺序返回记忆树中所有任务级摘要（L3 目标摘要节点的自然语言摘要）。对应"列出你做了哪些任务"类问题，直接给出按时间排列的任务清单，Agent 无需展开任何节点即可回答 task-list 类问题。

\textbf{（2）task\_lookup(query, max\_matches=8)：}在任务级摘要中使用嵌入向量匹配与 query 最相关的 top-k 个任务，返回每个匹配的摘要文本和对应的时间范围。对应"描述了某个任务的细节"类问题。

\textbf{（3）object\_lookup(object\_name, max\_matches=8)：}在所有 L2 事件摘要中查找提及指定物体名称（支持子串匹配）的事件，并按时间排序。内部先做字符串子串匹配，再对候选集进行嵌入向量精排。对应"某物体出现在哪些任务中"类问题。关键防护：对长度 $\le 3$ 的短物体名（如 "cd"），要求精确匹配而非子串匹配，防止其漂移到 "credit card" 等噪声匹配。

\textbf{（4）date\_lookup(query, max\_matches=8)：}解析 query 中的日期或相对天数表达式（如 "21 days ago"、"November 23, 2023"），在树中定位对应时间范围内的事件摘要。返回格式包含日期、事件摘要和语义匹配分数。对应"某天做了什么"类问题。

\textbf{（5）event\_date\_lookup(query, max\_matches=12)：}针对"how many days ago did you do X"类问题的专用工具。解析 query 中的事件描述和目标日期，找到所有匹配事件并计算距目标日期的天数。返回格式为 "score=X; days\_ago=Y; date=Z: summary"。

\textbf{（6）temporal\_neighbor(target\_task, direction, max\_candidates=6)：}给定一个目标任务描述，在树中定位对应事件，然后沿时间方向（before/after）返回相邻的任务候选。内部通过 task\_lookup 定位目标 → 沿时间相邻边或树结构找到前后邻居。对应"X 任务之前/之后做了什么"类问题。

\textbf{（7）action\_lookup(query, max\_matches=6)：}针对低层动作问题，将原始 L2 事件级动作映射回其所属的任务级摘要。对应"当你执行某个具体动作时正在完成什么任务"类问题。当低层动作已通过结构化工具给出高置信答案时，Agent 直接使用 \texttt{answer()} 返回，不再进入自由工具循环，这种"结构化直答"机制可有效减少低难度问题上的无效 token 消耗。

\subsection{缓存体系与执行护栏}

长历史评测的稳定性不仅依赖于检索算法本身，还依赖于执行层面的系统性控制。GRAF-Mem 在评测管线上构建了三层执行护栏。

\textbf{（1）预处理缓存体系。}每个长历史（如由 50 个 episode 拼接而成）在正式 QA 评测前，先通过预处理流程批量调用大语言模型摘要器生成完整的层级摘要树（递归概括至单一根节点），并将结果序列化为缓存文件。正式评测阶段启用 \texttt{require-cache} 强制只读模式，确保历史只能从缓存中加载，禁止在 QA 过程中触发任何在线摘要生成。评测前还需进行缓存完整性审计，确认所有必需文件均已生成且可读。

\textbf{（2）checkpoint 与容错。}长时间运行的评测支持 checkpoint / resume 机制：每个样本完成后立即写入 JSON 输出文件，中断后重新启动时自动跳过已完成的样本。对于瞬时 API 错误（网络超时、限流等），单样本支持最多 2 次自动重试。

\textbf{（3）token 与时间上限。}设置单样本 prompt token 上限和平均 prompt token 上限，防止个别样本因反复搜索导致 token 消耗失控，进而拖慢整轮实验。单样本运行时间上限用于防止死循环或 API 长时间无响应造成阻塞。

\section{TEACh 长历史问答实验}

\subsection{实验设置}

本章的实验在 TEACh 数据集上进行。TEACh 是面向家庭环境的具身交互数据集，包含按时间展开的任务执行轨迹与自然语言问答，适合评估长历史条件下的检索与问答能力，因此本文选用它作为实验平台。本文用 $|h|$ 表示一段长历史中拼接的基础 episode 数，其中 episode 指一个完整的基础情景片段。实验使用 $|h|=50$ 和 $|h|=100$ 两种历史长度设置；每种设置包含 10 段长历史（每段由 50 或 100 个 episode 拼接而成），每段配有 10 种不同类型的问答对，共计 100 题。TEACh 数据集的详细背景和问题类型分布见第二章 2.7.1 节。

\textbf{基线}：采用原论文中 H-EMV 在 TEACh 上报告的 full multimodal ICL=1 结果作为参考对照。该结果来自文献报告而非本文复现实验，因此本文仅用于数量级比较。

\textbf{GRAF-Mem 配置}：使用图增强检索标准配置，包含事件关联图（默认启用时间相邻、共享物体、共位置、相似动作四类边；因果边作为可选项，不纳入主实验）和全部七种结构化工具。Agent 模型为 Gemini-2.5-Pro（通过 API 调用），嵌入模型为 all-mpnet-base-v2。所有评测历史均使用预处理缓存，正式 QA 中强制 \texttt{require-cache}，无在线摘要成本。

\textbf{评测指标}：
\begin{itemize}
\item \textbf{S\textsubscript{c}（Semantic Correctness，完全正确率）}：回答在语义上完全命中标准答案的比例，由独立于 Agent 的大语言模型裁判（LLM Judge，Gemini-2.5-Pro）判定为 CORRECT 的比例
\item \textbf{S\textsubscript{p}（Semantic Partial，至少部分正确率）}：回答至少抓住关键事实的比例，即 Judge 判定为 CORRECT 或 PARTIAL 的比例
\item \textbf{T（每题平均 prompt token）}：每个 QA 的平均 prompt token 消耗量，衡量检索和问答的上下文成本
\end{itemize}

\subsection{主实验结果}

\begin{table}[htbp]
\centering
\caption{TEACh 长历史问答主结果}
\label{tab:3-1}
\begin{tabular}{|l|c|c|c|c|}
\hline
实验设置 & $|h|$ & S\textsubscript{c} & S\textsubscript{p} & T（每题平均 prompt token） \\
\hline
H-EMV baseline (ICL=1) & 50 & 37.0\% & 68.0\% & 10.2K \\
\hline
\textbf{GRAF-Mem} & \textbf{50} & \textbf{48.0\%} & \textbf{72.0\%} & \textbf{2.06K} \\
\hline
H-EMV baseline (ICL=1) & 100 & 34.0\% & 62.0\% & 10.4K \\
\hline
\textbf{GRAF-Mem} & \textbf{100} & \textbf{49.0\%} & \textbf{75.0\%} & \textbf{2.29K} \\
\hline
\end{tabular}
\end{table}

\begin{table}[htbp]
\centering
\caption{GRAF-Mem 相对参考对照的优势量化}
\label{tab:3-2}
\begin{tabular}{|l|c|c|c|}
\hline
对比项 & $\Delta$S\textsubscript{c} & $\Delta$S\textsubscript{p} & $\Delta$T \\
\hline
$|h|=50$ GRAF-Mem vs Baseline & +11.0 pp & +4.0 pp & \textbf{$-$79.8\%} \\
\hline
$|h|=100$ GRAF-Mem vs Baseline & +15.0 pp & +13.0 pp & \textbf{$-$78.0\%} \\
\hline
\end{tabular}
\end{table}

从表 3-1 和表 3-2 可以得出三个关键观察：

\textbf{（1）语义正确率显著提升。}在 $|h|=50$ 设置下，S\textsubscript{c} 从 37.0\% 提升至 48.0\%（+11.0 个百分点）；在更长的 $|h|=100$ 设置下，提升更为显著：S\textsubscript{c} 从 34.0\% 提升至 49.0\%（+15.0 个百分点），S\textsubscript{p} 从 62.0\% 提升至 75.0\%（+13.0 个百分点）。GRAF-Mem 不仅在整体上优于参考对照，而且随着历史长度的增加，优势反而扩大。参考对照在 $|h|$ 从 50 增至 100 时 S\textsubscript{c} 下降了 3 个百分点（37\%→34\%），而 GRAF-Mem 基本持平（48\%→49\%）。

\textbf{（2）token 成本大幅降低。}在两个设置下，GRAF-Mem 的平均 prompt token 消耗均约为 2.1K–2.3K，仅为参考对照的 20\% 左右，降低约 78\%–80\%。H-EMV 的 10.2K–10.4K prompt tokens 中包含大量将完整历史直接拼接成超长提示词后的展开内容和在线搜索上下文，而 GRAF-Mem 的预处理缓存体系将这部分成本完全剥离。更值得注意的是，token 成本的降低并不以牺牲准确率为代价，恰恰相反，在更低的 token 预算下取得了更高的语义正确率。

\textbf{（3）对历史长度变化的敏感度较低。}当历史长度从 50 个 episode 翻倍到 100 个 episode 时，GRAF-Mem 的 S\textsubscript{c} 仅从 48.0\% 变为 49.0\%（+1.0 pp），T 仅从 2.06K 增至 2.29K（+11.2\%）。而参考对照在相同翻倍下 S\textsubscript{c} 下降了 3.0 pp，T 反而微增（10.2K→10.4K，+2.0\%）。这说明图增强检索和结构化工具在一定程度上\textbf{抵消了历史增长带来的检索难度}：事件图的横向扩展和工具的精确命中使得 Agent 不需要在更大的树中盲目搜索更多节点。

\subsection{消融实验}

本节进一步评估不同边类型对检索性能的贡献。所有设置均沿用与主实验一致的 100 题评测集、相同的 Agent、相同的结构化工具和相同的缓存管线，仅改变图中激活的边类型。\texttt{full graph} 对应时间相邻 + 共享物体 + 共位置 + 相似动作四类边同时启用；\texttt{temporal-only} 仅启用时间相邻边；\texttt{object-only} 仅启用共享物体边。

\begin{figure}[htbp]
    \centering
    \includegraphics[width=\textwidth]{图3-4 图边类型消融实验结果柱状图.png}
    \caption{图边类型消融实验结果}
    \label{fig:chapter3-ablation}
\end{figure}

\medskip
{\small\textbf{注：}此处 baseline 指仅保留原始纯语义检索、不启用事件图扩展的设置；与 3.5.2 节中的 H-EMV 参考对照不同。}

从图 3-4 可以看出：

\textbf{（1）完整图结构带来最好的整体性能。}full graph 相比 baseline 的 S\textsubscript{c} 从 36.0\% 提升到 42.0\%（+6.0 个百分点），S\textsubscript{p} 从 56.0\% 提升到 65.0\%（+9.0 个百分点）。与此同时，平均 prompt token 仅从 4.57K 小幅增加到 4.76K，说明图扩展带来的上下文开销是可控的，而准确率收益更为明显。

\textbf{（2）时间边主要增强局部时序推理。}temporal-only 虽然不如 full graph 全面，但已经将 S\textsubscript{c} / S\textsubscript{p} 提升到 38.0\% / 61.0\%，且 T 下降到 4.41K。结合题型分布来看，它在 \texttt{right after}、\texttt{right before} 这类时序题上的提升最明显，说明时间相邻边对"前后发生了什么"类问题具有直接贡献。

\textbf{（3）物体边主要增强实体关联召回。}object-only 将 S\textsubscript{c} / S\textsubscript{p} 提升到 39.0\% / 59.0\%，在三种单一边设置中表现最好。它在 \texttt{object yes/no}、\texttt{where else} 这类物体关联题上的提升最明显，说明共享物体边更擅长把同一实体跨事件的证据串联起来。

综合来看，时间边提供顺序约束，物体边提供实体桥接，而 full graph 将二者与空间和动作关联一起整合，因此在总体上取得最佳结果。这也说明第三章的图增强检索并不是依赖单一关系，而是依赖多种横向语义信号的互补。

\subsection{结果分析与机制归因}

GRAF-Mem 相对参考对照的双重优势，即语义正确率更高、token 成本更低，不宜归因于单一因素，而是以下几个层面的协同效应。

\textbf{（1）事件图扩展了检索覆盖面。}H-EMV 的纯文本语义搜索在面对需要横向关联的问题时容易遗漏相关证据。例如，用户问"这个杯子后来出现在哪些任务中"，语义搜索可能返回包含"cup"或"mug"的事件，但可能遗漏那些虽涉及同一杯子实例、却没有直接写出"cup"一词的事件（例如摘要只写了"I placed the dirty object in the sink"，而 "dirty object" 指向的就是那个杯子）。事件图的共享物体边通过 Jaccard 加权连接共享同一物体实例 ID 的所有事件，不依赖文本措辞；只要视觉场景图构建时正确地分配了同一实例 ID，就能建立连接。同样，时间相邻边确保了时序推理问题的候选事件不被遗漏，共位置边为空间推理提供了一致的关联路径。

\textbf{（2）结构化工具降低了搜索噪声。}参考对照中，Agent 需要在将完整历史直接拼接成超长提示词的文本中寻找答案，或通过通用的 expand/search 在树结构中自由探索。当问题涉及特定日期或特定物体时，Agent 可能反复展开不相关的分支、进行多余的语义搜索，导致 token 消耗虚高且可能被噪声信息干扰。GRAF-Mem 的结构化工具（特别是 date\_lookup、event\_date\_lookup、object\_lookup）将这种盲目搜索替换为精确的结构化查找，Agent 调用一次工具即可获得按相关性排序、包含时间/日期/分数的结构化返回，无需在树结构中逐层展开和折叠。实际运行中，task\_lookup 和 date\_lookup 的引入通常能够减少 Agent 的自由搜索轮次。

\textbf{（3）缓存和上限控制降低了评测阶段的附加开销。}参考对照的 10.2K–10.4K prompt tokens 不仅包含 Agent 主动搜索引入的上下文，还可能包含评测过程中因历史状态复制与在线摘要而产生的附加开销。GRAF-Mem 的预处理缓存 + require-cache 机制在正式 QA 开始前就确保了所有 episode 的完整层级摘要已生成并固化为缓存文件，评测中不再触发额外的大语言模型摘要调用。此外，低层动作结构化直答（action\_lookup 返回高置信结果时直接 \texttt{answer()} 而不进入自由工具循环）和短对象护栏共同抑制了常见的无效 token 消耗模式。

综合来看，GRAF-Mem 的优势可以被概括为：\textbf{不是在让 Agent 看更多的上下文，而是在让 Agent 看到更干净、更相关的证据集合}。事件图提供横向联想能力以提升召回率，结构化工具提供精确命中能力以提升精确率，缓存和上限控制确保这两者在不产生额外在线摘要开销的前提下运作。

\section{本章小结}

本章针对层级记忆树纵向检索在横向语义关联上的不足（RQ1），提出了事件关联图构建和图增强检索流水线。事件关联图从 L2 事件节点自动派生，包含五种语义边（时间相邻、共享物体、共位置、相似动作、因果），每种边均有明确的构建算法和密度控制策略。图增强检索通过"语义种子召回→图邻居扩展→联合重排"三阶段流水线将横向关联信号嵌入检索链路，重排公式中 $\alpha \cdot \text{语义项} + \beta \cdot \text{图贡献项} + \gamma \cdot \text{深度奖励项}$ 的三项加权设计确保了语义相关性始终主导排序。结构化检索工具将记忆树中的任务、物体、日期等结构信息暴露为七种可调用工具，配合预处理缓存和执行护栏，将 Agent 在长历史上的搜索行为从自由探索转变为结构化命中。

在 TEACh 数据集上的实验结果表明，GRAF-Mem 在 $|h|=100$ 设置下达到 S\textsubscript{c}=49.0\%、S\textsubscript{p}=75.0\%、T=2.29K，相比参考对照 S\textsubscript{c} 提升 15 个百分点、token 成本降低 78\%。更值得注意的是，在 $|h|$ 从 50 翻倍到 100 时 GRAF-Mem 的 S\textsubscript{c} 基本持平，说明图增强检索在一定程度上抵消了历史增长带来的检索难度。消融实验进一步表明，时间边和物体边分别对时序题与物体关联题贡献最明显，而 full graph 通过多种边类型互补取得最佳整体性能。机制分析表明，这一双重优势来自事件图的横向扩展、结构化工具的精确命中以及缓存与上限控制的协同。
